\documentclass[11pt]{article}
\usepackage[T1]{fontenc}
\usepackage{lmodern}
\usepackage[a4paper,margin=27mm]{geometry}
\usepackage{amsmath,amssymb,amsthm,mathtools}
\usepackage{microtype,booktabs,longtable,array}
\usepackage{xcolor,listings,fancyhdr,needspace}
\usepackage{hyperref}
\definecolor{linkink}{HTML}{294E47}
\hypersetup{colorlinks=true,linkcolor=linkink,citecolor=linkink,urlcolor=linkink,
 pdftitle={Exact computation and cubic-order growth of OEIS A281434},
 pdfauthor={Research note prepared for Vladimir Reshetnikov}}
\pagestyle{fancy}
\fancyhf{}
\fancyhead[L]{\small Exact computation of OEIS A281434}
\fancyhead[R]{\small September 2026}
\fancyfoot[C]{\thepage}
\setlength{\headheight}{14pt}
\setlength{\parskip}{4pt}
\setlength{\parindent}{0pt}
\setlength{\emergencystretch}{2em}
\numberwithin{equation}{section}
\newtheorem{theorem}{Theorem}[section]
\newtheorem{proposition}[theorem]{Proposition}
\newtheorem{lemma}[theorem]{Lemma}
\newtheorem{corollary}[theorem]{Corollary}
\theoremstyle{definition}
\newtheorem{definition}[theorem]{Definition}
\theoremstyle{remark}
\newtheorem{remark}[theorem]{Remark}
\DeclareMathOperator{\supp}{supp}
\newcommand{\Z}{\mathbb Z}
\newcommand{\N}{\mathbb N}
\newcommand{\fall}[2]{#1^{\underline{#2}}}
\newcommand{\rise}[2]{#1^{\overline{#2}}}
\newcommand{\D}{\frac{d}{dx}}
\newcommand{\epsn}{\varepsilon_n}
\lstset{language=Python,basicstyle=\small\ttfamily,columns=fullflexible,
 keepspaces=true,showstringspaces=false,breaklines=true,frame=single,
 framerule=0.3pt,rulecolor=\color{black!25},xleftmargin=5pt,xrightmargin=5pt,
 aboveskip=10pt,belowskip=10pt}

\title{\textbf{Exact computation and cubic-order growth\\of the expanded derivatives of $x^{x^x}$}\\[5pt]
\large OEIS A281434}
\author{Research note prepared for Vladimir Reshetnikov}
\date{September 20, 2026}

\begin{document}
\maketitle
\begin{abstract}
Let $a(n)$ be the number of nonzero monomials in the fully expanded $n$th
 derivative of $x^{x^x}$, after collecting identical terms. We replace repeated
 symbolic differentiation by a seven-transition integer recurrence on three
 exponents. It computes the prefix through $n$ in $O(n^4)$ integer arithmetic
 operations and $O(n^3)$ live integers. A faster implementation performs the
 recurrence modulo a machine-size integer and verifies every possible zero
 by an independent exact coefficient formula; this method is deterministic,
 not a probabilistic zero test. We prove the explicit bounds
\[
 \frac{7n^3+39n^2+26n+3\epsn(n-1)}{24}
 \ \le a(n)\le\ \frac{2n^3+3n^2+4n}{3},\qquad n\ge1,
\]
where $\epsn$ is $1$ for odd $n$ and $0$ for even $n$. In particular,
$a(n)=\Theta(n^3)$. The upper bound counts an explicit lattice envelope;
the lower bound uses nonvanishing leading coefficients and high-multiplicity
roots of coefficient polynomials. We also identify an unavoidable family of
logarithm-free cancellations and an infinite family with positive logarithm
exponent. The stronger statement $a(n)\sim \tfrac23n^3$, and the still stronger
linear-deficit conjecture, are \emph{not proved here}. Exact code, tests,
terms through $n=100$, and selected computations through $n=200$ accompany
this article.
\end{abstract}

\clearpage
\tableofcontents
\clearpage

\section{The problem and the scope of the results}

The sequence A281434 is defined in the OEIS as the number of terms in the fully
expanded $n$th derivative of $x^{x^x}$, with offset $0$~\cite{oeis281434}.
Exponentiation is right-associated. Thus the function in question is
\[
 f(x)=x^{(x^x)}=\exp\bigl(x^x\log x\bigr),\qquad x>0,
\]
not $(x^x)^x=x^{x^2}$. Its initial term counts are
\[
1,3,12,30,64,113,188,285,415,577,780,1017,1312,\ldots.
\]
The positivity restriction supplies an unambiguous real logarithm; all
subsequent coefficient manipulations are exact polynomial identities.

The phrase ``fully expanded'' must include collecting like terms and deleting
zero coefficients. Counting product-rule branches is a different problem.
Even recording every monomial that could be reached by differentiation gives
an overcount: distinct contributions can cancel exactly.

For $n\ge1$, define
\begin{equation}\label{eq:B}
 B(n)=\frac{2n^3+3n^2+4n}{3},\qquad
 \Delta(n)=B(n)-a(n).
\end{equation}
The polynomial $B(n)$ is A037237$(n-1)$~\cite{oeis037237}. The related OEIS
entry A352697 records the differences $\Delta(n)$ and conjectures both
$\Delta(n)=O(n)$ and infinitely many zero differences~\cite{oeis352697}.
These observations are useful motivation, but neither is needed by our
algorithms or proofs.

\begin{theorem}[Main results]\label{thm:main}
There is a deterministic exact algorithm for A281434 requiring $O(n^4)$
integer arithmetic operations and $O(n^3)$ stored integers to compute
$a(0),\ldots,a(n)$. There is also a machine-modular implementation for a
single $a(n)$, with exact verification of every candidate cancellation.
For $n\ge1$,
\begin{equation}\label{eq:mainbounds}
 L(n):=\frac{7n^3+39n^2+26n+3\epsn(n-1)}{24}
 \le a(n)\le B(n),
 \quad \epsn=\frac{1-(-1)^n}{2}.
\end{equation}
Consequently
\begin{equation}\label{eq:theta}
 a(n)=\Theta(n^3),\qquad
 \lim_{n\to\infty}\frac{\log a(n)}{\log n}=3.
\end{equation}
\end{theorem}

The distinction between a growth \emph{order} and an asymptotic
\emph{equivalent} is important. Equation~\eqref{eq:theta} is unconditional.
An equivalent $a(n)\sim\tfrac23n^3$ would additionally require
$\Delta(n)=o(n^3)$; the present bounds do not establish that. The computations
are consistent with a much smaller deficit, but this numerical observation
will not be substituted for a nonvanishing theorem.

\section{A canonical differential-polynomial representation}

\subsection{Three auxiliary variables}
Put
\[
 u=x^x,\qquad v=x^{-1},\qquad z=\log x,\qquad A=z(z+1).
\]
The three elementary derivatives and the logarithmic derivative of $f$ are
\begin{equation}\label{eq:basicderivatives}
 u'=u(z+1),\qquad v'=-v^2,\qquad z'=v,
 \qquad \frac{f'}{f}=u(A+v).
\end{equation}
There is a unique polynomial $P_n\in\Z[u,v,z]$ such that
\begin{equation}\label{eq:Pdefinition}
 f^{(n)}(x)=f(x)P_n(x^x,x^{-1},\log x).
\end{equation}
Write
\begin{equation}\label{eq:coeffdefinition}
 P_n(u,v,z)=\sum_{k,j,\ell\ge0}c_n(k,j,\ell)u^kv^jz^\ell.
\end{equation}
The corresponding expanded term of the original derivative is
\[
 c_n(k,j,\ell)\,x^{x^x+kx-j}(\log x)^\ell.
\]
Therefore
\begin{equation}\label{eq:count}
 a(n)=\#\{(k,j,\ell):c_n(k,j,\ell)\ne0\}.
\end{equation}
In particular, a negative coefficient contributes one term, not minus one.

\subsection{Why the representation is unambiguous}

Suppose a finite linear combination of $u^kv^jz^\ell$, after substitution,
vanishes for all sufficiently large positive $x$. Group it by $k$ and choose
the largest $k$ having a nonzero Laurent-polynomial coefficient in $x$ and
$\log x$. Division by $x^{kx}$ makes every smaller-$k$ group exponentially
small relative to that coefficient. Within the largest group, choose the
smallest $j$, and then the largest $\ell$ with nonzero coefficient. Its
asymptotic magnitude cannot be canceled by the remaining Laurent-logarithmic
terms. This is a contradiction. Hence the displayed monomials are linearly
independent in the required finite sense.

This argument is also a safeguard against a representation-dependent term
count: we are counting a specified canonical basis, not the summands in an
arbitrarily factored expression.

\subsection{The recurrence and its proof}

\begin{proposition}[Differential recurrence]\label{prop:recurrence}
Starting from $P_0=1$, the polynomials satisfy
\begin{equation}\label{eq:PRec}
 P_{n+1}=u(z^2+z+v)P_n+u(z+1)\partial_uP_n
          -v^2\partial_vP_n+v\partial_zP_n.
\end{equation}
\end{proposition}
\begin{proof}
Differentiate~\eqref{eq:Pdefinition}, apply the product rule, and substitute
\eqref{eq:basicderivatives}. The result is exactly~\eqref{eq:PRec}.
The base case is $f^{(0)}=f$. Induction proves the identity at every order,
and also proves integrality of every coefficient.
\end{proof}

For one coefficient $c=c_n(k,j,\ell)$ the seven contributions are as follows.
All contributions with the same destination must be added before zero
coefficients are removed.
\begin{center}
\begin{tabular}{lll}
\toprule
Destination at order $n+1$ & Contribution & Source in the derivative\\
\midrule
$(k+1,j,\ell+2)$ & $c$ & $f'/f$: $uz^2$\\
$(k+1,j,\ell+1)$ & $c$ & $f'/f$: $uz$\\
$(k+1,j+1,\ell)$ & $c$ & $f'/f$: $uv$\\
$(k,j,\ell+1)$ & $kc$ & differentiation of $u^k$\\
$(k,j,\ell)$ & $kc$ & differentiation of $u^k$\\
$(k,j+1,\ell)$ & $-jc$ & differentiation of $v^j$\\
$(k,j+1,\ell-1)$ & $\ell c$ & differentiation of $z^\ell$\\
\bottomrule
\end{tabular}
\end{center}
Transitions with a zero multiplier are omitted. In particular, no negative
logarithm exponent is ever introduced.

\subsection{The first cancellation}

The first two polynomials have the compact forms
\begin{align}
 P_1&=u(A+v),\label{eq:P1}\\
 P_2&=u^2(A+v)^2+
 u\bigl((z+1)A+(3z+2)v-v^2\bigr).\label{eq:P2}
\end{align}
Expansion gives $a(1)=3$ and $a(2)=12$. At the next order, the coefficient
polynomial of $uv^2$ is
\[
 -(z+1)+\bigl(\partial_z-1\bigr)(3z+2)=-4z.
\]
Its constant term vanishes. The candidate envelope has $31$ cells at order
$3$, but the derivative has $30$ terms. A Boolean reachability calculation
would fail at this very early example.

\section{The support envelope and the cubic upper bound}

\begin{definition}
For $n\ge1$, $1\le k\le n$, and $0\le j\le n$, put
\begin{align}
 m_{n,k,j}&=\max(0,k-j),\label{eq:lo}\\
 U_{n,k,j}&=
 \begin{cases}
 n+k,&j=0,\\
 \min(n+k-j-1,\,2n-2j),&j\ge1.
 \end{cases}\label{eq:hi}
\end{align}
Let $\mathcal E_n$ consist of the integer triples with these ranges and
$m_{n,k,j}\le\ell\le U_{n,k,j}$.
\end{definition}

\begin{proposition}[Support containment]\label{prop:envelope}
For $n\ge1$, $\supp P_n\subseteq\mathcal E_n$.
\end{proposition}
\begin{proof}
Consider a nonzero differentiation path, and let $a,b,c,d,e,r,s$ count its
uses of the seven transitions in the displayed table, in that order. Then
\[
 k=a+b+c,\quad j=c+r+s,\quad
 \ell=2a+b+d-s,\quad
 n=a+b+c+d+e+r+s.
\]
In particular,
\begin{align*}
 n-j&=a+b+d+e,\\
 \ell-(k-j)&=a+d+r\ge0,\\
 2(n-j)-\ell&=b+d+2e+s\ge0,\\
 n+k-j-\ell&=b+c+e+s\ge0.
\end{align*}
If $j>0$, at least one transition of type $c$ or $s$ must occur: the transition
of type $r$ has multiplier $-j$ and cannot create the first positive $j$.
Thus $n+k-j-\ell\ge1$ when $j>0$. The remaining ranges follow immediately
from the transition table and the first differentiation. Every coefficient
is a sum over such paths, so cancellation can remove cells but cannot add a
cell outside these constraints.
\end{proof}

\begin{proposition}[Size of the envelope]\label{prop:envelopesize}
For $n\ge1$, $\#\mathcal E_n=B(n)$.
\end{proposition}
\begin{proof}
For $j=0$ there are $n+1$ logarithm exponents for each of $n$ values of $k$,
giving $n(n+1)$ cells. For $j\ge1$, set $d=n-j$, so $0\le d\le n-1$.
The number of cells in this $j$-slice is
\begin{align*}
 \sum_{k=1}^n\bigl(\min(d+k,2d+1)-\max(0,k-n+d)\bigr)
 &=n(2d+1)-d(d+1).
\end{align*}
Indeed, replacing each minimum by $2d+1$ overcounts by
$d(d+1)/2$, and the sum of the subtracted maxima is another $d(d+1)/2$.
Consequently
\begin{align*}
 \#\mathcal E_n
 &=n(n+1)+\sum_{d=0}^{n-1}\bigl(n(2d+1)-d(d+1)\bigr)\\
 &=n(n+1)+n^3-\frac{(n-1)n(n+1)}3\\
 &=\frac{2n^3+3n^2+4n}{3}.
\end{align*}
\end{proof}

This proves the upper bound in~\eqref{eq:mainbounds}. More precisely,
$\Delta(n)$ is exactly the number of zero coefficients \emph{inside}
$\mathcal E_n$. The polynomial $B(n)$ is not being guessed by interpolation:
it is the exact count of a proved finite set of candidate monomials.

\section{An independent exact formula for every coefficient}

The recurrence is already an exact algorithm. An independent formula is
valuable for two other reasons: it certifies possible cancellations after a
modular computation, and it proves that every $(k,j)$ row is nonzero.

\subsection{Notation for Stirling numbers and auxiliary polynomials}

We use signed Stirling numbers $s(n,k)$ of the first kind and Stirling
numbers $S(n,k)$ of the second kind, with the conventions
\[
 \fall{y}{n}=\prod_{i=0}^{n-1}(y-i)=\sum_k s(n,k)y^k,
 \qquad y^n=\sum_k S(n,k)\fall{y}{k}.
\]
These conventions agree with DLMF \S26.8~\cite{dlmf}. In particular,
$s(n,k)$ has sign $(-1)^{n-k}$ and is nonzero for $1\le k\le n$;
$S(n,k)>0$ on the same range.

For nonnegative $d,p,q$, define the noncentral Stirling numbers
\begin{equation}\label{eq:Tdef}
 T_d(p;q)=\sum_{r=0}^d\binom dr q^{d-r}S(r,p).
\end{equation}
They can be computed without rational arithmetic:
\begin{equation}\label{eq:Trec}
 T_0(p;q)=\mathbf1_{p=0},\qquad
 T_{d+1}(p;q)=(p+q)T_d(p;q)+T_d(p-1;q).
\end{equation}
Indices $p<0$ or $p>d$ have value zero. The useful boundary values are
\[
 T_d(0;q)=q^d,\qquad T_d(d;q)=1,\qquad
 T_d(p;0)=S(d,p),\qquad T_d(p;1)=S(d+1,p+1),
\]
where $0^0=1$. Equivalently,
\begin{equation}\label{eq:Tmoment}
 e^{-w}\sum_{m\ge0}(m+q)^d\frac{w^m}{m!}
 =\sum_{p=0}^dT_d(p;q)w^p.
\end{equation}
To verify this identity, expand $(m+q)^d$, express powers of $m$ in falling
factorials, and sum $\sum_m\fall{m}{p}w^m/m!=w^pe^w$.

Define
\begin{equation}\label{eq:Qdef}
 Q_{j,r}(\beta)=\fall{\beta}{j}\,\rise{(\beta+1)}{r}
 =\prod_{i=0}^{j-1}(\beta-i)\prod_{i=1}^{r}(\beta+i).
\end{equation}
These integer polynomials are easy to build by multiplying by linear
factors. Notice that $\deg Q_{j,r}=j+r$.

\subsection{A parameterized derivative identity}

For parameters $\alpha,\beta$, Taylor expansion with the increment $xt$ gives
\begin{equation}\label{eq:taylorparam}
 \frac{(x+xt)^{\alpha(x+xt)+\beta}}{x^{\alpha x+\beta}}
 =(1+t)^\beta
 \exp\!\left(\alpha x\,[tz+(1+t)\log(1+t)]\right).
\end{equation}
Comparing coefficients of $t^n$ and then of $(\alpha x)^d$ yields
\begin{equation}\label{eq:parameteridentity}
 \frac{d^n}{dx^n}x^{\alpha x+\beta}
 =x^{\alpha x+\beta-n}\sum_{d=0}^n(\alpha x)^d
 C_{n,d}(z,\beta),
\end{equation}
where, with $j=n-d$,
\begin{equation}\label{eq:Cformula}
 C_{n,d}(z,\beta)
 =\sum_{s=0}^d\binom ns\frac{z^s}{(d-s)!}
 \partial_\beta^{\,d-s}Q_{j,d-s}(\beta).
\end{equation}
For completeness, choose $s$ factors $tz$ from the $d$th power of the bracket
in~\eqref{eq:taylorparam}. The remaining factors contribute
$(1+t)^{d-s}\log^{d-s}(1+t)$. Write the logarithm power as
$\partial_\beta^{d-s}$ acting on $(1+t)^{\beta+d-s}$ and use
\[
 [t^{n-s}](1+t)^{\beta+d-s}
 =\frac{\fall{(\beta+d-s)}{n-s}}{(n-s)!}
 =\frac{Q_{j,d-s}(\beta)}{(n-s)!}.
\]
The factorial prefactors simplify to $\binom ns/(d-s)!$, proving the formula.

\subsection{The coefficient oracle}

\begin{theorem}[Exact coefficient formula]\label{thm:oracle}
Suppose $(k,j,\ell)\in\mathcal E_n$, and put
\[
 d=n-j,\qquad h=d+k-\ell,
\]
\[
 q_- =\max(0,k-\ell,k-d),\qquad q_+=\min(k,h,j).
\]
Then
\begin{equation}\label{eq:oracle}
\boxed{\displaystyle
 c_n(k,j,\ell)=\sum_{q=q_-}^{q_+}
 \binom n{\ell-k+q}\binom hq\,
 T_d(k-q;q)\,[\beta^h]Q_{j,h-q}(\beta).}
\end{equation}
An empty sum is zero. Outside the envelope, the coefficient is zero.
\end{theorem}
\begin{proof}
Expand the outer exponential:
\[
 f(x)=\sum_{m\ge0}\frac{x^{mx}z^m}{m!},\qquad
 x^{mx}z^m=\left.\partial_\beta^m x^{mx+\beta}\right|_{\beta=0}.
\]
Apply~\eqref{eq:parameteridentity} with $\alpha=m$. For any polynomial
$C(\beta)$,
\[
 \left.\partial_\beta^m\bigl(e^{\beta z}C(\beta)\bigr)\right|_{\beta=0}
 =\sum_{q=0}^m\frac{m!}{(m-q)!}z^{m-q}[\beta^q]C(\beta).
\]
After division by $f=e^{uz}$, the sum over $m$ occurring for fixed $d,q$ is
\begin{align*}
 e^{-uz}\sum_{m\ge q}\frac{m^d u^m z^{m-q}}{(m-q)!}
 &=u^q e^{-uz}\sum_{r\ge0}(r+q)^d\frac{(uz)^r}{r!}\\
 &=u^q\sum_{p=0}^dT_d(p;q)(uz)^p.
\end{align*}
To obtain $u^kz^\ell$, take $p=k-q$ and take
$s=\ell-k+q$ from~\eqref{eq:Cformula}. Set $r=d-s=h-q$.
Finally,
\[
 [\beta^q]\frac{\partial_\beta^rQ_{j,r}(\beta)}{r!}
 =\binom{q+r}{q}[\beta^{q+r}]Q_{j,r}(\beta),
\]
and $q+r=h$. This proves~\eqref{eq:oracle}; the stated bounds on $q$ enforce
$p\in[0,d]$, $s\in[0,d]$, and $q\le j$.

The infinite series can be read as locally normally convergent analytic
series for positive $x$ near any fixed point, or as formal exponential
identities followed by extraction of the finitely many relevant coefficients.
\end{proof}

\subsection{Three boundary checks}

Useful independently checkable consequences are
\begin{align}
 [v^0]P_n&=(z+1)^n\sum_{k=1}^nS(n,k)(uz)^k,
                 &&n\ge1,\label{eq:vzero}\\
 [u^n]P_n&=(z(z+1)+v)^n,\label{eq:umax}\\
 [v^n]P_n&=\fall{u}{n}.\label{eq:vmax}
\end{align}
The first follows either from~\eqref{eq:oracle} or from setting $v=0$ in the
recurrence. The second follows because reaching $u$-degree $n$ requires every
step to differentiate the outer factor $f$. For the third, the coefficient
of the next highest $v$-power is obtained by multiplication by $u-n$,
starting from $1$.

\section{A rigorous cubic lower bound}

An upper bound alone does not establish cubic growth. We now show that a
positive cubic number of the candidate cells must survive, without
assuming that cancellations are rare.

\subsection{Every row has its advertised highest coefficient}

Let
\[
 R_{n,k,j}(z)=[u^kv^j]P_n(u,v,z).
\]

\begin{lemma}[Nonzero leading coefficient]\label{lem:leading}
For every $1\le k\le n$ and $0\le j\le n$, the polynomial $R_{n,k,j}$ is
nonzero and has degree exactly $U_{n,k,j}$.
\end{lemma}
\begin{proof}
For $j=0$, equation~\eqref{eq:vzero} gives
\[
 R_{n,k,0}(z)=S(n,k)z^k(z+1)^n.
\]
Now let $j\ge1$ and put $d=n-j$. If $k\le d$, the proposed highest exponent
is $\ell=d+k-1$, so $h=1$ in~\eqref{eq:oracle}. Only $q=0,1$ contribute.
The coefficient is
\begin{equation}\label{eq:leadlowk}
 (-1)^{j-1}(j-1)!\left(
 \binom n{d-1}S(d,k)+\binom nd S(d+1,k)\right),
\end{equation}
which is nonzero. Here we used
$[\beta]Q_{j,0}=[\beta]Q_{j,1}=(-1)^{j-1}(j-1)!$.

If $k\ge d+1$, the highest exponent is $\ell=2d$, and $h=k-d$.
The only possible summand is $q=h$, giving
\begin{equation}\label{eq:leadhighk}
 [z^{2d}]R_{n,k,j}(z)=\binom nd s(j,k-d).
\end{equation}
It is nonzero because $1\le k-d\le j$. The support envelope already rules
out higher exponents, completing the proof.
\end{proof}

\subsection{A high-multiplicity root}

\begin{lemma}[Divisibility]\label{lem:divisibility}
The row polynomial is divisible by
\begin{equation}\label{eq:divisibility}
 z^{\max(k-j,0)}(z+1)^{M_{n,k,j}},
 \qquad M_{n,k,j}=\max(n-2j,k-j,0).
\end{equation}
\end{lemma}
\begin{proof}
Extracting $u^kv^j$ from~\eqref{eq:PRec} gives
\begin{align}\label{eq:rowrec}
 R_{n+1,k,j}={}&z(z+1)R_{n,k-1,j}+R_{n,k-1,j-1}
              +k(z+1)R_{n,k,j}\notag\\
 &+\bigl(\partial_z-(j-1)\bigr)R_{n,k,j-1}.
\end{align}
Rows with out-of-range indices are zero. Begin with $R_{0,0,0}=1$.
For the factor $z+1$, the first and third terms increase the inherited
multiplicity by one. The second term has inherited multiplicity at least
$\max(n-2j+2,k-j,0)$. The derivative in the last term can reduce its inherited
multiplicity by at most one; before differentiation it is at least
$\max(n-2j+2,k-j+1,0)$. Each is sufficient for the target exponent
$\max(n+1-2j,k-j,0)$. This proves the $(z+1)$ assertion by induction.
The same argument with the factor $z$ proves the exponent $\max(k-j,0)$.
Since $z$ and $z+1$ are relatively prime, both divisibilities hold together.
\end{proof}

\begin{lemma}[Multiplicity forces many terms]\label{lem:sparse}
A nonzero polynomial with a nonzero root of multiplicity at least $M$ has
at least $M+1$ nonzero coefficients.
\end{lemma}
\begin{proof}
Suppose $p(z)=\sum_{i=1}^t b_i z^{e_i}$, with distinct nonnegative integers
$e_i$ and $b_i\ne0$. If a root $\alpha\ne0$ has multiplicity at least $t$,
then $(z\partial_z)^r p(\alpha)=0$ for $0\le r<t$. Thus
\[
 \sum_{i=1}^t e_i^r b_i\alpha^{e_i}=0,\qquad 0\le r<t.
\]
The matrix $(e_i^r)$ is an invertible Vandermonde matrix, contradicting
$b_i\alpha^{e_i}\ne0$. Therefore the multiplicity is at most $t-1$.
\end{proof}

\subsection{Summing the mandatory terms}

By Lemmas~\ref{lem:leading}--\ref{lem:sparse}, each $(k,j)$ row contributes at
least $1+M_{n,k,j}$ nonzero monomials. Rows are disjoint in the canonical
basis, so
\begin{equation}\label{eq:lowersum}
 a(n)\ge\sum_{k=1}^n\sum_{j=0}^n
 \bigl(1+\max(n-2j,k-j,0)\bigr).
\end{equation}
For fixed $j$, the sum of the maxima over $k$ is
\begin{equation}\label{eq:Mslice}
 \sum_{k=1}^n M_{n,k,j}=
 \begin{cases}
 n(n-2j)+\dfrac{j(j+1)}2,&0\le j\le\lfloor n/2\rfloor,\\[5pt]
 \dfrac{(n-j)(n-j+1)}2,&j>\lfloor n/2\rfloor.
 \end{cases}
\end{equation}
In the first case, start with $n-2j$ in each row and add the excess
$1+\cdots+j$. In the second case, only $k>j$ contributes.
Summing these elementary progressions, and adding the $n(n+1)$ ones
in~\eqref{eq:lowersum}, gives exactly $L(n)$ in~\eqref{eq:mainbounds}.
This completes the proof of the growth bounds in Theorem~\ref{thm:main}.

In particular,
\begin{equation}\label{eq:limbounds}
 \frac7{24}\le\liminf_{n\to\infty}\frac{a(n)}{n^3}
 \le\limsup_{n\to\infty}\frac{a(n)}{n^3}\le\frac23.
\end{equation}
These constants are proved bounds, not estimates of an experimentally fitted
leading coefficient.

\section{Exact cancellation families}

\subsection{Logarithm-free coefficients and derivatives of a falling factorial}

Let $F_n(y)=\fall{y}{n}$. In~\eqref{eq:oracle}, setting $\ell=0$ forces
$q=k$, so for $d=n-j$ and $1\le k\le j$,
\begin{equation}\label{eq:logfree}
\boxed{\displaystyle
 c_n(k,j,0)=\frac{k^d}{d!\,k!}F_n^{(d+k)}(d).}
\end{equation}
Although written using factorial quotients, this expression is an integer;
the implementation instead uses integer polynomial coefficients and binomial
coefficients. Formula~\eqref{eq:logfree} reduces every logarithm-free
cancellation to an exact derivative evaluation of a single elementary
polynomial.

\begin{proposition}[Central cancellations at every odd order]\label{prop:odd}
For odd $n\ge1$,
\begin{equation}\label{eq:odddeficit}
 \Delta(n)\ge\left\lfloor\frac{n+1}{4}\right\rfloor.
\end{equation}
\end{proposition}
\begin{proof}
Write $n=2d+1$ and set $j=d+1$. Then
\[
 F_n(d+\beta)=\beta\prod_{r=1}^d(\beta^2-r^2)
\]
is odd. For each $1\le k\le d+1$ with $k\equiv d\pmod2$, its coefficient of
$\beta^{d+k}$ vanishes. Formula~\eqref{eq:logfree} gives a missing monomial
$(k,d+1,0)$. There are $\lfloor(d+1)/2\rfloor=\lfloor(n+1)/4\rfloor$ such
values of $k$.
\end{proof}

Thus $a(n)=B(n)$ is impossible at every odd $n\ge3$. This does not resolve
whether equality occurs infinitely often at even orders. Nor does it
classify all logarithm-free cancellations: additional integer derivative
zeros occur away from the central symmetry.

\subsection{An infinite family with positive logarithm exponent}

It is tempting to guess that all cancellations are logarithm-free.
The following family disproves that guess.

Since $f=e^g$ with $g=uz$, the $u^{n-1}$ part of $f^{(n)}/f$ consists of the
partition type having one block of size two and all other blocks of size one.
Equivalently, induction in~\eqref{eq:PRec} gives
\begin{equation}\label{eq:almosttop}
 [u^{n-1}]P_n=\binom n2(A+v)^{n-2}
 \bigl((z+1)A+(3z+2)v-v^2\bigr),\qquad n\ge2.
\end{equation}
For $2\le j\le n-1$, put $m=n-j-1$ and
\[
 a=\binom{n-2}{j},\qquad b=\binom{n-2}{j-1},\qquad
 c=\binom{n-2}{j-2}.
\]
The coefficient of $v^j$, after removing the common factor $\binom n2$, is
\begin{equation}\label{eq:quadraticrow}
 z^m(z+1)^m\left[-cz^2+(a+3b-c)z+(a+2b)\right].
\end{equation}
Its coefficient of $z^{2m+1}$ is
\begin{equation}\label{eq:nexttotop}
 a+3b-(m+1)c.
\end{equation}
Using
\[
 \frac ac=\frac{m(m+1)}{j(j-1)},\qquad
 \frac bc=\frac{m+1}{j-1},
\]
this vanishes exactly when $m=j^2-4j$ in the stated range.

\begin{proposition}[Non-logarithm-free cancellation family]\label{prop:family}
For every integer $t\ge2$, let
\begin{equation}\label{eq:family}
 n=t^2+t-1,\quad k=n-1,\quad j=t+2,\quad \ell=2t^2-7.
\end{equation}
Then $(k,j,\ell)\in\mathcal E_n$, $\ell>0$, and
$c_n(k,j,\ell)=0$.
\end{proposition}
\begin{proof}
Substitute $j=t+2$ into $m=j^2-4j$: it gives $m=t^2-4$ and
$n=m+j+1=t^2+t-1$. Then $\ell=2m+1$ is positive, is inside the envelope,
and its coefficient vanishes by~\eqref{eq:nexttotop}. The next higher
coefficient is $-\binom n2c\ne0$, so this is a genuine internal missing term.
\end{proof}

The first corresponding orders are $5,11,19,29,41,55,\ldots$, with missing
triples
\[
 (4,4,1),\ (10,5,11),\ (18,6,25),\ (28,7,43),\ldots.
\]
These are additional to the logarithm-free central cancellations. Since all
orders in~\eqref{eq:family} are odd, this proves
$\Delta(n)\ge\lfloor(n+1)/4\rfloor+1$ along that infinite subsequence.

\section{Algorithms and complexity}

\subsection{The dependency-free exact recurrence}

The most direct implementation stores a dictionary from triples to integer
coefficients. At each differentiation it applies the seven transitions,
combines equal keys, and removes zeros. A short reference version is given
in Appendix~\ref{app:code}. It uses no symbolic-algebra system and no
numerical approximation.

The supplied production version packs a triple into one integer key. For a
requested maximum order $N$, take
\[
 b=N+2,\qquad c=2N+3,\qquad \operatorname{key}(k,j,\ell)=(kb+j)c+\ell.
\]
Incrementing $k,j,\ell$ corresponds to adding $bc,c,1$, respectively. The
chosen bases exceed every possible exponent, so no carries can identify
different monomials. The packed representation reduces tuple allocation
and dictionary overhead; it does not change the mathematics.

\begin{proposition}[Arithmetic complexity]\label{prop:complexity}
The exact recurrence computes the entire prefix through $N$ in $O(N^4)$
integer arithmetic operations and $O(N^3)$ live integer coefficients.
Each coefficient has $O(N\log N)$ bits.
\end{proposition}
\begin{proof}
At order $n$ there are at most $B(n)=O(n^3)$ terms, and each contributes to
at most seven destinations. Summing through $N$ gives $O(N^4)$ operations.
Only the current and next dictionaries are needed, giving $O(N^3)$ storage.

Let $H_n$ be the sum of the absolute values of the coefficients of $P_n$.
The absolute sum of all multipliers leaving a term is
$3+2k+j+\ell\le3+5n$. Therefore
\[
 H_{n+1}\le(5n+3)H_n,\qquad H_0=1,
 \qquad H_n\le\prod_{r=0}^{n-1}(5r+3)\le5^n n!.
\]
Every individual coefficient satisfies this bound, so its bit length is
$O(n\log n)$.
\end{proof}

This is a unit-cost \emph{integer-arithmetic} bound, not a claim that
arbitrarily large integers are machine words. With elementary small-integer
multiplication, a conservative bit bound is $O(N^5\log^2N)$ time and
$O(N^4\log N)$ coefficient-storage bits. Hashing and object overhead affect
practical memory. We do not claim the algorithm is asymptotically optimal.

\subsection{A deterministic modular sieve with exact refinement}

Most coefficients need never be constructed as large integers. Fix any
integer modulus $M\ge2$, and run the same recurrence in $\Z/M\Z$.
If a residue is nonzero, the original integer coefficient is certainly
nonzero. Only zero residues need further attention.

The exact single-term algorithm is:
\begin{enumerate}
\item Compute $P_n\bmod M$ by the seven-transition recurrence.
\item Enumerate the triples in $\mathcal E_n$ whose residues are zero.
\item Evaluate each such coefficient with~\eqref{eq:oracle}, using integer
      arithmetic, and count only the coefficients that are actually zero.
\item Return $B(n)$ minus that exact zero count.
\end{enumerate}
A zero residue is \emph{never} accepted as proof of a zero coefficient.
Consequently this algorithm is correct for every admissible modulus,
including small or composite moduli. No randomness, primality assumption,
or conjecture about support is required. Tests deliberately use moduli
$2$, $7$, and $12$ to force many false modular zeros.

The supplied NumPy implementation uses signed 64-bit arrays and the modulus
$1\,000\,000\,007$ by default. At order $m$, the array has shape
$(m+1,m+1,2m+1)$. The next array has shape
$(m+2,m+2,2m+3)$, and each transition becomes a shifted slice addition.
All entries are reduced modulo $M$ after each step. The implementation checks
\[
 (10n+10)M<2^{63}
\]
before starting; this conservative condition bounds the intermediate
products and sums before modular reduction. Overflow would invalidate the
certificate, so the check is part of correctness rather than merely an
optimization detail.

The dense modular pass takes $O(n^4)$ fixed-width operations and $O(n^3)$
fixed-width words. Let $Z_M(n)$ denote the number of zero residues inside
the envelope. The oracle tables $Q_{j,r}$ and $T_d(p;q)$, and Pascal's
triangle for binomial coefficients, are built lazily. For a single target
$n$, filling all potentially needed tables takes $O(n^3)$ integer arithmetic
operations and storage; each candidate coefficient then requires at most
$O(n)$ integer products and additions. Thus a conservative oracle bound is
\[
 O(n^3+nZ_M(n))\subseteq O(n^4).
\]
In practice the lazy tables are much smaller when the modular pass leaves
few candidates. Unlike the reference recurrence, the oracle may multiply
two large integers, so its unit-arithmetic bound should again not be read
as a bit-operation bound.

\subsection{Which implementation to use}

For a whole prefix, the dependency-free recurrence has the simplest predictable
cost and automatically supplies every intermediate count. For an isolated
larger index, the modular method usually avoids nearly all large-integer
coefficient work. Running a new modular pass separately for every prefix
index would discard this advantage and incur repeated work.

Neither method differentiates a transcendental expression in a computer
algebra system. Both return the exact OEIS term count, including signed
cancellations, rather than an expression-size estimate.

\section{Computational validation and results}

\subsection{Checks actually performed}

The accompanying files record ten passing test groups. The tests compare the
reference recurrence with all $43$ published A281434 values through $n=42$,
and compare the generated prefix with all $80$ published A352697 deficits
through $n=80$~\cite{oeis281434,oeis352697}. Every candidate coefficient
through $n=12$ is checked against the independent formula~\eqref{eq:oracle}.
The two methods are also compared through $n=12$ for each of four moduli,
including the deliberately small and composite choices described above.

Additional tests verify the envelope, its exact cardinality, the lower-bound
sum, the $(z+1)$ multiplicities, the two leading-coefficient formulas,
the odd central cancellations, and initial instances of the infinite
non-logarithm-free family. Such finite tests check the implementations and
help detect algebraic mistakes; the universal statements are justified by
the proofs, not by the tests.

The integer recurrence was run through $n=100$. The modular implementation
was separately checked against it at $n=20,40,60,80,100$. Selected larger
computations used the modular pass and exact verification of all candidate
zeros. The archive contains the complete generated prefix, the tests, and
the selected missing triples.

\subsection{Selected exact values}

\begin{center}
\begin{tabular}{rrrr}
\toprule
$n$ & $a(n)$ & $B(n)$ & $\Delta(n)$\\
\midrule
0 & 1 & \multicolumn{1}{c}{--} & \multicolumn{1}{c}{--}\\
1 & 3 & 3 & 0\\
2 & 12 & 12 & 0\\
3 & 30 & 31 & 1\\
5 & 113 & 115 & 2\\
11 & 1\,017 & 1\,023 & 6\\
20 & 5\,760 & 5\,760 & 0\\
40 & 44\,320 & 44\,320 & 0\\
60 & 147\,680 & 147\,680 & 0\\
80 & 347\,838 & 347\,840 & 2\\
100 & 676\,800 & 676\,800 & 0\\
101 & 697\,178 & 697\,203 & 25\\
150 & 2\,272\,700 & 2\,272\,700 & 0\\
151 & 2\,318\,265 & 2\,318\,303 & 38\\
200 & 5\,373\,600 & 5\,373\,600 & 0\\
\bottomrule
\end{tabular}
\end{center}
For example, the two missing triples at $n=80$ are
$(33,36,0)$ and $(42,45,0)$. At $n=101$, all $25$ missing terms are supplied
by the central logarithm-free symmetry in Proposition~\ref{prop:odd}.
These are exact statements about the recorded finite computations.

\subsection{Measured performance}

On the execution environment used for this report, the packed integer
recurrence generated the full prefix $0\le n\le100$ in approximately
$26.16$ seconds. Selected single-index modular runs were:
\begin{center}
\begin{tabular}{rrrr}
\toprule
$n$ & Modular zero candidates & Actual zeros & Elapsed seconds\\
\midrule
40 & 0 & 0 & 0.021\\
60 & 0 & 0 & 0.096\\
80 & 2 & 2 & 0.342\\
100 & 0 & 0 & 0.820\\
101 & 25 & 25 & 0.698\\
150 & 0 & 0 & 4.970\\
151 & 38 & 38 & 5.460\\
200 & 0 & 0 & 17.769\\
\bottomrule
\end{tabular}
\end{center}
These are individual measurements, not a controlled benchmark suite or a
promise of timings on other hardware. In particular, computing one target
by the modular method and computing an entire prefix by the integer method
are different workloads. The exact timings, interpreter information, and
modulus are saved in \texttt{data/benchmarks.csv} and
\texttt{data/environment.json}.

The fact that every zero residue in this small selection was a genuine zero
is not assumed by the algorithm. The small-modulus tests exercise the
opposite situation and confirm that false modular zeros are rejected by
the exact oracle.

\section{What is and is not established asymptotically}

The unconditional conclusion is cubic-order growth, with the fully explicit
bounds~\eqref{eq:mainbounds}. Consequently, $a(n)=n^{3+o(1)}$ in the
logarithmic-exponent sense. The lower bound remains valid even if many more
cancellations occur at orders far beyond the computed range.

There are three progressively stronger claims:
\begin{align*}
 \text{Proved here:}\qquad &a(n)=\Theta(n^3).\\
 \text{Not proved here:}\qquad &a(n)\sim\frac23n^3.\\
 \text{Stronger OEIS conjecture:}\qquad
 &a(n)=\frac23n^3+n^2+O(n).
\end{align*}
The last statement is equivalent to $\Delta(n)=O(n)$, given the exact
formula for $B(n)$, and implies the middle statement. Neither follows from
agreement with $B(n)$ at many computed indices.

A proof of the middle statement needs only a subcubic bound on the number
of holes; it need not classify every cancellation. In particular, the
entire plane $\ell=0$ has only $O(n^2)$ candidate cells. Thus even a complete
failure to classify the logarithm-free derivative zeros would not by itself
prevent proving the leading equivalent. The remaining issue is a uniform
bound on cancellations with $\ell>0$ across the three-dimensional envelope.
Formula~\eqref{eq:oracle} supplies an explicit integer expression for those
coefficients, and the family in Proposition~\ref{prop:family} shows why a
blanket nonvanishing assertion would be false.

For the stronger linear-deficit conjecture, even the logarithm-free plane
requires more delicate control. Its central symmetry already forces
$\Delta(n)$ to be at least a positive linear function along the odd orders,
so an $o(n)$ bound is impossible. The appropriate target is an upper bound
of matching order, not an assertion that the derivative eventually fills
its whole envelope.

The computation problem is therefore resolved by unconditional exact
algorithms, and the polynomial growth exponent is established. Determining
the exact leading asymptotic constant remains a separate nonvanishing
problem in this analysis.

\clearpage
\appendix
\section{A compact exact implementation}\label{app:code}

The following version emphasizes transparency. The archive's
\texttt{code/a281434.py} additionally provides packed keys, the coefficient
oracle, modular arrays, input validation, and a command-line interface.

\begin{lstlisting}
from collections import defaultdict

def a281434_prefix(N):
    if not isinstance(N, int) or isinstance(N, bool) or N < 0:
        raise ValueError("N must be a nonnegative integer")
    p = {(0, 0, 0): 1}
    values = [1]
    for _ in range(N):
        q = defaultdict(int)
        for (k, j, ell), c in p.items():
            q[k+1, j, ell+2] += c
            q[k+1, j, ell+1] += c
            q[k+1, j+1, ell] += c
            if k:
                q[k, j, ell+1] += k*c
                q[k, j, ell] += k*c
            if j:
                q[k, j+1, ell] -= j*c
            if ell:
                q[k, j+1, ell-1] += ell*c
        p = {key: c for key, c in q.items() if c}
        values.append(len(p))
    return values
\end{lstlisting}

The only potentially negative transition is the one with multiplier $-j$.
Deleting it, taking absolute values before accumulation, or replacing the
dictionary by a set changes the sequence.

\clearpage
\section{Reproducing the supplied computations}

Python 3.9 or later is sufficient for the dependency-free method. The modular
method additionally requires NumPy. From the archive's \texttt{code}
directory, the commands below produce a single term, inspect missing
monomials, create a prefix, and run the tests.
\begin{lstlisting}[language={}]
python a281434.py 100 --method modular
python a281434.py 80 --method modular --holes
python a281434.py 100 --prefix
python -m unittest -v test_a281434.py
python run_experiments.py --reference
python run_experiments.py 20 40 60 80 100 101 150 151 200
\end{lstlisting}
The first command returns \texttt{676800}. The second returns
\texttt{347838} and the two missing triples listed in the article.
The full prefix and related deficits are in \texttt{data/b281434.txt} and
\texttt{data/b352697.txt}. These are generated data files accompanying this
report, not claims that the OEIS has incorporated the additional terms.

The LaTeX source builds with
\begin{lstlisting}[language={}]
pdflatex -interaction=nonstopmode -halt-on-error article.tex
pdflatex -interaction=nonstopmode -halt-on-error article.tex
pdflatex -interaction=nonstopmode -halt-on-error article.tex
\end{lstlisting}
Repeated runs resolve the table of contents and cross-references. The bundled
\texttt{build.sh} and \texttt{build.ps1} perform these commands.

\Needspace{15\baselineskip}
\begin{thebibliography}{9}
\bibitem{oeis281434}
Vladimir Reshetnikov,
\emph{A281434: Number of terms in the fully expanded $n$-th derivative of
$x^{(x^x)}$}, The On-Line Encyclopedia of Integer Sequences, entry created
October 5, 2017.
\url{https://oeis.org/search?q=id:A281434&fmt=text}.
Accessed September 20, 2026.

\bibitem{oeis352697}
Vladimir Reshetnikov,
\emph{A352697: $a(n)=\mathrm{A037237}(n-1)-\mathrm{A281434}(n)$},
The On-Line Encyclopedia of Integer Sequences, March 29, 2022.
\url{https://oeis.org/A352697}.
Accessed September 20, 2026.

\bibitem{oeis037237}
N. J. A. Sloane and contributors,
\emph{A037237: Expansion of $(3+x^2)/(1-x)^4$},
The On-Line Encyclopedia of Integer Sequences.
\url{https://oeis.org/A037237}.
Accessed September 20, 2026.

\bibitem{dlmf}
NIST Digital Library of Mathematical Functions,
\emph{\S26.8: Set Partitions---Stirling Numbers}.
\url{https://dlmf.nist.gov/26.8}.
Accessed September 20, 2026.
\end{thebibliography}

\end{document}
